\documentclass[11pt]{article}
\usepackage[T1]{fontenc}
\usepackage{lmodern,amsmath,amssymb,amsthm,booktabs}
\usepackage[margin=1in]{geometry}
\usepackage[colorlinks=true,linkcolor=blue,urlcolor=blue,citecolor=blue]{hyperref}
\usepackage{microtype}
\newtheorem{theorem}{Theorem}
\newtheorem{lemma}[theorem]{Lemma}
\newtheorem{corollary}[theorem]{Corollary}
\newcommand{\OPT}{\operatorname{OPT}}
\newcommand{\tri}[1]{\tau(#1)}
\title{Certified Shortcuts and Depth Bounds\\for Full Optimal Binary Search Trees\\
\large Research checkpoint: exact shortcuts and barriers to relaxations}
\author{OpenAI Codex Sydney\\
\small Problem selection and sponsorship: the repository owner}
\date{October 18, 2026}
\begin{document}
\maketitle
\begin{abstract}
We investigate exact optimal binary search trees with ordered weights on
both keys and gaps. We give sufficient key and block contraction rules,
with construction cost $O(n+m^2+\sum_t k_t^2)$ for a surviving core of
$m$ keys and contracted blocks of sizes $k_t$. Complementary shortcuts
optimize minimum or bounded-excess unweighted path-length classes and
certify unrestricted optimality with residual-weight bounds. Successful
narrow-class attempts take $O(n\log n)$ work. A known split-height Monge
construction supplies an $O(n^{3/2})$ attempted route at every size, with a
separate certificate against taller trees.

We also implement the established short-interval/top-skeleton architecture
using a variable mass contour. A directly proved mixed-weight Fibonacci
bound guarantees that the contour contains every optimal subtree at a
chosen depth $d$. If $S_d$ intervals are admitted, exact full construction
takes $O(n\log(n+1)+\min_d\{S_d+n2^d\})$ arithmetic work, including
parameter selection. A cancellation argument bounds all short-table root
searches by $3S_d$. On a positive family with a localized light block, the
variable contour gives $O(n^{1.59024})$ work while the specified uniform
width envelope requires $\Omega(n^{1.75})$ work. This is a comparison
within one architecture, not a new general subquadratic bound.

Exact enumeration, independent cubic recurrences, adversarial inputs, and
operation counts accompany all implementations. We preserve a failed
mass-bound transfer and document an ordered depth relaxation, its
$O(n\log n)$ oracle, and obstructions to local repairs. General worst-case
construction remains quadratic; novelty and broad practical superiority
are not established. The full subquadratic research objective remains open.
\end{abstract}

\section{Model, scope, and prior work}
Keys have fixed inorder positions $0,\ldots,n-1$ and nonnegative weights
$p_0,\ldots,p_{n-1}$. The $n+1$ gaps have weights $q_0,\ldots,q_n$.
A key at depth $d$ contributes $p(d+1)$; a gap at depth $d$ contributes $qd$.
Charging an extra comparison at each gap adds the tree-independent constant
$\sum q_i$.

For a half-open key interval $[i,j)$, write
\[
 W(i,j)=\sum_{r=i}^{j-1}p_r+\sum_{r=i}^{j}q_r,\qquad C(i,i)=0.
\]
The exact recurrence is
\[
 C(i,j)=W(i,j)+\min_{i\le r<j}\{C(i,r)+C(r+1,j)\}.
\]
Knuth's monotone-root optimization computes the table in $O(n^2)$
time~\cite{knuth}. The central obstacle here is the number of states, not
merely the cost per state. We retain both kinds of weights throughout.

Larmore gives subquadratic exact construction under restrictions on long
runs of very low frequency items, together with an approximation
tradeoff~\cite{larmore}. Karpinski, Larmore, and Rytter combine exact small
intervals with a bounded top-level construction~\cite{klr}.
These results are essential context: an improvement on a restricted family
must not be described as the first subquadratic result.
This checkpoint does not make a comprehensive priority claim. The conditions
below were derived in this investigation; a targeted literature search did
not establish whether equivalent rules already exist.

All running-time bounds below count arithmetic operations and comparisons.
They are not bit-complexity bounds for unbounded integer weights.

\section{A sharper single-key contraction}
A key is a \emph{cherry} when its two children are external nodes. Replacing
a cherry with one gap is exact: if the key and its gaps have weights
$p,x,y$, their total is $s=p+x+y$. A cherry rooted at depth $d$ contributes
$s(d+1)$, whereas the replacement gap contributes $sd$.

\begin{theorem}[Single-key rule]
Let a key of weight $p$ have adjacent gap weights $x,y$.
If its immediate predecessor and successor keys have weights $a,b$ and
\[
 \min(a,b)\ \ge\ 2p+x+y,
\]
there is an optimal tree in which this key is a cherry.
At the left endpoint, it suffices that the successor weight is at least
$p+x$; at the right endpoint, it suffices that the predecessor weight is
at least $p+y$. Consequently contraction to a gap of weight $s=p+x+y$
satisfies $\OPT(\text{original})=\OPT(\text{contracted})+s$.
\end{theorem}

\begin{proof}
Consider the key at depth $d$ in any tree. If its left subtree is nonempty,
the immediate predecessor has depth $d+l$, $l\ge1$. If its right subtree is
nonempty, the immediate successor has depth $d+r$, $r\ge1$.
Set the corresponding distance to zero for an empty internal side.

When $l\ge1$, move the predecessor to the chosen key's position, replacing
the predecessor's old position by its left child. Insert the chosen key as
a cherry at the leftmost gap of the former right subtree. Inorder order is
preserved. All other outcomes either retain their depths or become shallower.
The predecessor saves $la$, and the cost change is at most
\begin{equation}\label{eq:left}
 -la+(r+1)p+(r-l+1)x+y.
\end{equation}
The mirrored successor exchange has change at most
\begin{equation}\label{eq:right}
 -rb+(l+1)p+x+(l-r+1)y.
\end{equation}

If both sides are nonempty and $l\ge r$, substituting $a\ge2p+x+y$
in~\eqref{eq:left} gives
\[
 (r+1-2l)(p+x)+(1-l)y\le0.
\]
Use~\eqref{eq:right} when $r\ge l$. If only the left side is nonempty,
$r=0$ and $a\ge p+y$ makes~\eqref{eq:left} nonpositive.
The other case is symmetric. The stated interior condition implies these
one-sided conditions. At a global endpoint only the appropriate one-sided
case is possible. With both sides empty the key is already a cherry.

Applying the exchange to an optimum proves the existence claim.
Contraction subtracts exactly $s$. Conversely, every contracted tree
expands to a valid original tree at additional cost $s$, proving equality.
\end{proof}

The weaker condition $\min(a,b)\ge p+x+y$ is false.
For $p=(24,14,24)$ and $q=(1,0,9,0)$, the middle key has local mass 23
and neighbors of weight 24. Contracting it forces cost 142, while the
unrestricted optimum costs 130.

Maintain a doubly linked list of remaining keys and weighted gaps.
Initially queue every key; after a contraction, only its predecessor and
successor need rechecking. At most $3n$ queue entries are processed.
Saving the contracted tree at the new gap permits reconstruction.
Thus this pass takes $O(n)$ operations and space, and a Knuth solve on the
remaining $m$ keys gives total time $O(n+m^2)$. No confluence or
minimum-core-size property is claimed.

\section{Whole-block contraction}
Let $B=[i,j)$ be a nonempty proper interval, including its endpoint gaps
$q_i,q_j$. Put $W=W(i,j)$ and $C=C(i,j)$, with costs measured from a
stand-alone root of depth zero.

\begin{theorem}[Block rule]\label{thm:block}
If both boundary keys exist and their weights satisfy
\[
 p_{i-1}\ge W+C,\qquad p_j\ge W+C,
\]
there is an optimum in which all of $B$ forms one subtree.
For a prefix or suffix block, it suffices that the sole boundary key
has weight at least $C$. Contracting the block to one gap of weight $W$
then gives
\[
 \OPT(\text{original})=\OPT(\text{contracted})+C.
\]
\end{theorem}

\begin{proof}
In any tree, let $v$ be the lowest common ancestor of gaps $q_i,q_j$,
at depth $d$. Its key lies in $B$, and every outcome of $B$ descends
from $v$. Hence the original contribution of $B$ is at least $(d+1)W$.

Within the subtree rooted at $v$, delete the keys of $B$ from the inward
fringes of its two children. On the left, a deleted key has no retained
right child and can be suppressed in favor of its left child; on the
right the process is symmetric. Let the retained outside trees be $A,D$.
All their retained outcomes are at depths no greater than before.
If $A$ is nonempty, its maximum is boundary key $i-1$; if $D$ is
nonempty, its minimum is boundary key $j$.

Suppose both are nonempty, and their boundary keys had original depths
$d+l,d+r$. If $l\ge r$, remove the maximum key from $A$, replacing its
position by its left child, and make that key the new root at $v$.
Its left child is the remainder of $A$ and its right child is $D$.
Put an optimal tree for $B$ at the leftmost external position of $D$.
This position has depth at most $d+r+1$.
No outside outcome other than the promoted boundary moves deeper,
and that boundary saves $l p_{i-1}$.
The increase in the contribution of $B$ is at most
\[
 C+(d+r+1)W-(d+1)W=C+rW.
\]
Since $l\ge r\ge1$ and $p_{i-1}\ge W+C$, the saving is at least
$l(W+C)\ge rW+C$. The mirrored construction applies for $r\ge l$.

If only one retained outside tree is nonempty, promote its inward boundary
key to $v$ and attach the optimal block as the other child.
The block increase is at most $C$, while that boundary saves at least
its own weight. The interior hypothesis suffices, and for a global prefix
or suffix the weaker stated hypothesis suffices. If both outside trees
are empty, $B$ already occupies the subtree at $v$ and can be optimized
in place.

The construction never increases cost. Apply it to an optimum.
A block at depth $e$ contributes $C+eW$, while its replacement gap
contributes $eW$. Contraction and expansion now prove the equality.
\end{proof}

A balanced $k$-key tree has height
$h_k=\lceil\log_2(k+1)\rceil$ and cost at most $h_kW$.
Thus the interior test
\[
 \min(p_{i-1},p_j)\ge(1+h_k)W
\]
and endpoint test $p_{\rm boundary}\ge h_kW$ certify a block without first
solving it. These are sufficient conditions, not characterizations.

\section{An adaptive linear selection pass}
Build the maximum Cartesian tree of the key-weight sequence, breaking ties
deterministically. Its subtree intervals form a laminar family: intervals
are disjoint or nested. Compute their sizes and masses by one traversal and
prefix sums. Test each proper interval with the balanced-tree bound above.
This inspects only $O(n)$ candidates. The restriction to Cartesian intervals
is an algorithmic choice, not a claim that all useful blocks occur there.

For a threshold $b$, select the maximal certified intervals whose sizes
are at most $b$. Different selected intervals have a surviving separator
key between them. Contracting one therefore leaves every other selected
interval's weights and boundary-key weights unchanged, so the block theorem
applies successively. Write their sizes as $k_1,\ldots,k_s$, and
$m=n-\sum_t k_t$. The exact number of nonempty interval-DP states is
\[
 S(b)=\tri m+\sum_t\tri{k_t},\qquad \tri u=\frac{u(u+1)}2.
\]

We choose a threshold minimizing $S(b)$ without running these DPs.
For every certified interval $I$ of size $k$, find its maximal proper
certified descendants $J$. Immediately before $I$ becomes eligible,
these descendants are its selected blocks. Its activation therefore
changes the core size and block-state total by
\[
 \Delta m=-\left(k-\sum_J |J|\right),\qquad
 \Delta S_{\rm blocks}=\tri k-\sum_J\tri{|J|}.
\]
These sums are computed bottom-up. Bucket activation events by integer
size $1,\ldots,n$ and scan the buckets, maintaining $m$ and the block-state
total. There are $O(n)$ events and buckets. Equal-size intervals are disjoint;
evaluate a threshold after processing its entire bucket.
The empty selection $b=0$ is included.

\begin{corollary}
The selected threshold and blocks are found in $O(n)$ arithmetic operations
and space. Solving the blocks and surviving core by Knuth's method and
expanding their saved trees takes
\[
 O\!\left(n+m^2+\sum_t k_t^2\right)
\]
time and $O(n+m^2+\max_t k_t^2)$ space when block tables are released
after each solve. The chosen state count is no larger than $n(n+1)/2$.
Worst-case time remains $O(n^2)$.
\end{corollary}

The minimization is only over these size thresholds in this candidate
hierarchy, not over every possible partition. It does not minimize
wall-clock time. Preprocessing, table allocation, root scans, and
reconstruction still count as work.

For perspective, a certified partition with
$m=O(n^{2/3})$ surviving keys and every $k_t=O(n^{1/3})$ gives
$O(n^{4/3})$ time, since $\sum k_t^2\le(\max k_t)n$.
This is a conditional structural consequence, not an unrestricted bound.

\section{Recursive construction with a bounded preprocessing depth}
The implementation \texttt{recursive\_contract.py} combines the proved
components. At each call it applies the single-key pass, compares the
order-free lower bound of Section~\ref{sec:unordered} with a feasible
bisection-tree cost, and returns immediately on equality. Otherwise it
selects blocks as above. It recurses only if their combined table-state
estimate is at most three quarters of the full estimate; small components
and unsuccessful decompositions fall back to Knuth's method.

Suppose a component has $m>8$ keys after the single-key pass and is
successfully decomposed. Every child component of size $s$ satisfies
$s(s+1)\le\tfrac34m(m+1)$, hence $s<0.92m$.
Child sizes sum to $m$. Consequently the recursion has logarithmic depth,
and the total number of keys processed on any recursion level is at most
$n$. Linear selection and reconstruction cost $O(n\log n)$ in total.
The heap-based equality checks cost $O(n\log^2 n)$ over all levels.
If the terminal DP sizes are $r_t$, the construction therefore takes
\[
 O\!\left(n\log^2 n+\sum_t r_t^2\right)
\]
arithmetic operations. Its worst case is still quadratic. This threshold
prevents repeated expensive scans after arbitrarily small reductions.
Correctness follows by contraction, exact equality certificates, and induction
on the strictly smaller recursive components; reconstruction restores the
original key indices and gap subtrees.

The recursive implementation passed 3,000 further random inputs through
$n=30$ against the cubic oracle and direct tree evaluation. A separate nested
fixture checks composition across multiple decomposition levels.
Matched Python timings use five inputs per family, alternate which solver
runs first, and include preprocessing and reconstruction for both solvers.
At $n=1024$, the median runtime ratio (recursive divided by Knuth) was
0.134 on the positive-weight heavy-separator family; its median terminal
DP count was 30,017 versus 524,800. On random mixed weights the ratio was
0.986 and on near-uniform weights 1.001, which support no meaningful gain.
At $n=512$, the respective random and near-uniform ratios were 1.131 and
1.093: preprocessing can make the method slower. These are exploratory
measurements on one M4 host, not production-language or cross-machine claims.

\newpage
\section{A certified shortcut through nearly complete shapes}
Let $I(T)$ be the sum of internal-node depths, with root depth zero.
The external path length is $I(T)+2n$, by induction on the two subtrees.
Consequently uniform key weight $a$ and gap weight $b$ give cost
\[
 (a+b)I(T)+(a+2b)n.
\]
For an arbitrary instance set $a=\min_i p_i$, $b=\min_j q_j$, and let
$H'$ be the order-free mixed-node Huffman bound of
Section~\ref{sec:unordered} on residual weights $p_i-a,q_j-b$.
These weights remain nonnegative. Write $I_*(n)$ for the minimum possible
unweighted internal path length, and let $C_*$ be the best weighted cost
among all trees with $I(T)=I_*(n)$.

\begin{samepage}
\begin{theorem}[Minimum-path-length certificate]
For $n\ge1$, if
\[
 C_*\ \le\ (a+b)(I_*(n)+1)+(a+2b)n+H',
\]
then a tree attaining $C_*$ is a globally optimal full OBST.
\end{theorem}
\end{samepage}
\begin{proof}
Split each outcome weight into its uniform baseline and residual. Every
tree outside the minimum-$I$ class has integer $I(T)\ge I_*(n)+1$, and its
residual cost is at least $H'$. Its full cost is therefore at least the
right-hand side. Within the class, $C_*$ is minimal by definition.
Equality is sufficient, including when $a+b=0$.
\end{proof}

\paragraph{Optimize the class without interval tables.}
Put $h=\lfloor\log_2(n+1)\rfloor$, $S=2^h$, $m=S-1$, and $r=n-m$.
Every minimum-$I$ tree has the complete upper skeleton with $m$ keys,
and expands exactly $r$ of its $S$ external slots into single-key cherries.
Indeed, a missing shallower slot could otherwise receive a deepest internal
leaf, strictly reducing $I$. Inorder relabeling preserves search-tree
feasibility in this unweighted shape argument. Thus
\[
 I_*(n)=(h-2)2^h+2+hr.
\]

Process skeleton slots $j=0,\ldots,S-1$ in inorder. State $(j,e)$ records
the number of processed slots and expansions, so the next input index is
$i=j+e$. An unexpanded slot charges $h q_i$; an expanded slot charges
$(h+1)(p_i+q_i+q_{i+1})$. Except after the last slot, also charge the next
skeleton key at comparison length
\[
 d_j=h-v_2(j+1),
\]
using $p_i$ after an unexpanded slot and $p_{i+1}$ after an expansion.
Here $v_2$ counts trailing zero bits. Exactly $r$ expansions are required,
so state $(j,e)$ is retained only when
\[
 \max(0,r-(S-j))\le e\le\min(r,j).
\]
With $s=\min(r,S-r)$, each stage has at most $s+1$ states and each state
has two outgoing choices. This computes $C_*$ and backpointers in
$O(n(1+s))$ time and space. A linear Cartesian-tree pass converts the
chosen comparison lengths into the labeled tree. No key or gap weights
are discarded in these transitions.

There is an exact early rejection test. Let $F$ be the cost obtained by
assigning independently sorted key and gap weights to the shallowest
depths of a minimum-$I$ shape. This is feasible for the order-free residual
relaxation, so $H'\le F-[(a+b)I_*+(a+2b)n]$. Therefore $C_*>F+a+b$
proves the certificate cannot succeed. Otherwise the sorted lists are
reused in a two-queue Huffman implementation: nonnegative weights and
nondecreasing internal weights make the sequence of merged weights
nondecreasing. This adds linear work after sorting. The queue optimization
is standard and is not claimed as novel.

The implementation attempts this DP only when
$s\le\lfloor\log_2 n\rfloor+1$. Including the residual bound,
preprocessing, and reconstruction, a successful attempt costs
$O(n\log n)$ arithmetic operations and space. A skipped or failed attempt
falls back to Knuth, or to the preceding recursive construction. Both
versions remain exact and have quadratic worst-case time. The test is a
certificate, not an assumption that near-uniform inputs have balanced optima.
Novelty of this shortcut has not been established.

\paragraph{Verification and limits.}
Minimum-$I$ optimization matched independent enumeration of all tree shapes
through eight keys on 360 weighted inputs. Both fallback versions matched
independent cubic DP and direct tree costs on 3,000 further mixed-weight
inputs through thirty keys; the unrestricted certificate succeeded on 768.
At $n=1,024$, one seeded input with every weight in $\{999,1000,1001\}$
used 2,048 shape states, 3,070 transitions, 10,511 sorting comparisons,
and 1,350 merge comparisons, certifying cost 19,479,590. The nonminimum-$I$ lower bound was
19,480,123. Knuth used 524,800 states and 1,050,139 root candidates.
These work units have different costs; the state ratio is not a timing claim.

The selected perturbation sweep certified amplitude 8 around 1,000 at
$n=256$, amplitude 4 at $n=512$, and amplitude 1 at $n=1,024$; larger tested
amplitudes at those sizes failed. These are seeded observations, not
universal thresholds. Ordinary random inputs failed the certificate.
Sizes 384, 768, and 1,000 deliberately skip the shortcut because their
minimum-$I$ shape space exceeds the logarithmic width budget.

Six matched trials per fixture used all six orderings of Knuth, recursive
contraction, and the combined solver, with identical inputs and correctness
checks outside timing. On Apple M4 with Python 3.9.6, median per-trial
combined/Knuth time ratios at $n=1,024$ were 0.03934 for near-uniform,
0.10478 for heavy separators, 1.05454 for random, 0.91929 for log-uniform,
and 1.03117 for gap spikes. All six near-uniform cases certified. The skipped
near-uniform controls at $n=768$ and $n=1,000$ gave ratios 1.17886 and 1.04896.
At $n=512$ the combined method was about 21\% slower than recursive
contraction alone on heavy separators. These exploratory timings include
sorting, lower bounds, and reconstruction; they do not establish broad
superiority. The final early-rejection and two-queue implementations were
also checked against the original heap certificate on all 3,000 inputs.

\paragraph{Optimize the baseline rather than fix it.}
The largest subtractable baseline need not give the strongest certificate.
Set $a_0=\min p_i$, $b_0=\min q_j$, and use
$a=a_0-\delta$, $b=b_0-\delta$ for an integer
$0\le\delta\le\lfloor(a_0+b_0)/2\rfloor$.
The residual weights remain nonnegative and $a+b\ge0$. Although one
baseline term may become negative, it is only an algebraic decomposition;
the input probabilities are unchanged. Let $H_\delta$ denote the residual
order-free optimum. The exclusion bound is
\[
 B(\delta)=(a_0+b_0-2\delta)(I_*+1)
              +(a_0+2b_0-3\delta)n+H_\delta.
\]
Every tree outside the minimum-$I$ class costs at least $B(\delta)$.
Thus any sampled price with $C_*\le B(\delta)$ certifies the same candidate.

For a fixed residual tree, increasing $\delta$ adds one unit to every
key and gap weight, changing its cost at rate $2I+3n$. Consequently
$B$ is a pointwise minimum of affine functions, hence concave. A minimizing
residual tree supplies a supergradient
\[
 g=2(I-I_*-1),\qquad
 B(x)\le B(\delta)+g(x-\delta).
\]
If $g>0$, smaller prices cannot improve the current value; if $g<0$,
larger prices cannot improve it. If $g=0$, the value is globally maximal.
Binary interval elimination, retaining the best sample, therefore finds
the maximum over allowed integer prices when run to completion.

The same linear oracle supplies $I$. Give every original outcome secondary
weight1, order queue items lexicographically by residual mass and number
of original outcomes, and sum both merge coordinates. The first gives
$H_\delta$ and the second gives $2I+3n$. This is the Huffman exchange
argument with an infinitesimal positive common perturbation, choosing
minimum $I$ among residual-cost ties. The weight order is independent of
$\delta$, so no additional sorting is needed. Each price costs $O(n)$.

The fixed-profile ceiling extends to all prices:
$B(\delta)\le F+a_0+b_0-2\delta$. Therefore an input rejected by the
existing profile test cannot be rescued by this search. The implementation
tries prices only after the old certificate fails while passing that test,
and permits at most $\lfloor\log_2n\rfloor+1$ extra oracle calls.
This preserves the $O(n\log n)$ attempted-shortcut bound. Reaching the cap
does not establish that the best price was found; every sampled certificate
is still valid, and failure invokes the exact fallback.

The lexicographic oracle matched 400 exhaustive shape-profile comparisons
through eight keys. Uncapped search matched exhaustive evaluation of all
allowed integer prices on 1,200 inputs. On 3,000 additional full-OBST
inputs through thirty keys, both capped fallback versions matched cubic
DP and direct tree costs; 87 previously unsuccessful unrestricted
certificates were recovered. A separate seeded sweep recovered the
$n=1,024$, amplitude 2 input with two additional oracle calls, and the
$n=2,048$, amplitude 1 input with three. Larger tested perturbations
usually exceeded the profile ceiling. This applies standard concave-price
optimization to the certificate; it is not a general subquadratic result.

In the seeded 2,048-key case, the candidate costs 43,032,446; the bound
rises from 43,031,347 at $\delta=0$ to 43,032,459 at $\delta=125$.
The three extra passes use 6,144 contractions and 6,141 pair comparisons,
in addition to 4,096 shape states, 6,142 transitions, 20,784 sorting
comparisons, and 2,752 original merge comparisons. Knuth uses 2,098,176
interval states and 4,198,710 root candidates on that input.

A separate six-trial comparison uses all orderings of Knuth, the old
hybrid, and the price-search hybrid. All six $n=1,024$, amplitude 2 cases
and all six $n=2,048$, amplitude 1 cases were newly certified. Median
per-trial price-search/Knuth ratios were respectively 0.04953 and 0.06128.
At $n=2,048$, amplitude 2 still failed, with ratio 0.99391. Ordinary random
and gap-spike controls at $n=1,024$ gave ratios 1.06492 and 1.12941.
The already-certified $n=1,024$, amplitude 1 control made no extra oracle
calls yet showed a price-search/old-hybrid ratio of 1.13564. This variability
limits interpretation of small timing differences; operation counts and
all raw observations are retained. The host was Apple M4 with Python 3.9.6.


\section{Exact bounded-excess construction by dyadic intervals}
The candidate class can be expanded beyond minimum internal path length.
This is necessary even on some mildly nonuniform inputs. For $n=255$, set
all gap weights to $16$ and all key weights to $16$, except keys $8j+1$
which have weight $17$. The unique minimum-$I$ tree costs $61{,}680$;
Knuth gives optimum $61{,}656$, attained
with $I=I_*+2$. The algorithm below constructs and certifies this optimum.
It retains both kinds of weight and never assumes that the restricted
optimum is globally optimal.

\paragraph{An exact representation.}
An ordered full binary tree is a tiling of $[0,1]$ by consecutive aligned
dyadic intervals: a gap at depth $d$ has width $2^{-d}$. A leaf starting
at $x$ must satisfy $x/2^{-d}\in\mathbb Z$. Conversely such a tiling is a
complete prefix code and determines an ordered full tree. Alignment, not
just the Kraft equality, enforces order. On a grid of height $H$, an
interior boundary at integer coordinate $Y$ corresponds to a key with
comparison length $H-\nu_2(Y)$. Thus gap and key costs are both local in
a left-to-right tiling construction.

\begin{theorem}[Bounded-excess construction]
For $n\ge1$, let $N=n+1$, $h=\lfloor\log_2N\rfloor$, $S=2^h$, $r=N-S$, and
$s=\min(r,S-r)$. For any nonnegative integer budget $k$, the minimum
full weighted cost among all trees with $I\le I_*+k$ can be found and
reconstructed in
\[
 O\bigl(n(s+k+1)(k+1)^2 2^k\bigr)
\]
arithmetic operations, using
$O\bigl(n(s+k+1)(k+1)2^k\bigr)$ space for backpointers.
\end{theorem}
\begin{proof}
Measure dyadic widths in units of $2^{-h}$ and put $w=2^{h-d}$.
For integral depth $d$, the potential
\[
 f(d)=d-h-2+2^{h+1-d}=d-h-2+2w
\]
is nonnegative, with zeros exactly at $h$ and $h+1$. The minimum external
path length is $hN+2r$. Since external path length equals $I+2n$ and
$\sum_i w_i=S$, every full tree satisfies
\[
 \sum_i f(d_i)=\sum_i d_i-(h+2)N+2S=I-I_*.
\]
Every depth in the budgeted class therefore has $f(d)\le k$. This implies
$d\le H=h+k+2$, and only $O(k+1)$ integral depths are possible.
For integral $d$, direct substitution gives
\[
 \tfrac12-\tfrac12 f(d)\le w\le1+f(d).
\]
Apply these inequalities to a prefix of $i$ leaves and to its remaining
suffix, each having potential at most $k$. The cumulative position $y$
must obey
\[
 \max\left(\tfrac{i-k}{2},i-r-k\right)
 \le y\le
 \min\left(i+k,\tfrac{S-r+i+k}{2}\right).
\]
This interval has width at most both $r+2k$ and $(S-r)/2+k$, hence
$O(s+k)$.

Set $K=2^{k+2}$. A state after $i$ leaves stores integral drift
$u=K(y-i)$ and spent potential $z=K\sum f(d)$. For depth $d$, let
$W=2^{H-d}$. Require $(iK+u)\bmod W=0$ and update
\[
 u'=u+W-K,\qquad z'=z+K(d-h)+2(W-K).
\]
Retain the position band and $0\le z'\le Kk$. Charge $q_i d$ and, for
$i>0$, charge $p_{i-1}(H-\nu_2(iK+u))$. Start at $(u,z)=(0,0)$ and
finish after $N$ leaves with $u=-rK$. Every retained complete path is an
aligned full tiling with excess at most $k$; every tree in the class
passes the same tests. The DP therefore optimizes exactly this class.
For fixed $u$, the congruence $z\equiv2u\pmod K$ leaves at most $k+1$
possible spent potentials. There are $O((s+k+1)2^k)$ positions per stage
and $O(k+1)$ choices per state, proving the bounds. The bands admit
direct array indexing; the Python prototype uses sparse dictionaries. Backpointers recover
gap depths; the boundary lengths reconstruct the tree in linear time.
\end{proof}

\paragraph{Perfect-size refinement.}
When $N=S$, use instead
$f_0(d)=d-h+\tfrac32(2^{h-d}-1)$. It is nonnegative on integer depths,
vanishes only at $h$, and also sums to $I-I_*$. The inequality
$|2^{h-d}-1|\le2f_0(d)$ and zero total drift imply $|y-i|\le k$.
Here $H=h+k+1$, $K=2^{k+1}$, and $z=2K\sum f_0$. Its increment is
$2K(d-h)+3(W-K)$, and $z\equiv3u\pmod{2K}$. The implementation uses this
smaller grid at perfect sizes, with the same asymptotic bound at $s=0$.

\paragraph{Certify against all other trees.}
Let $C_k$ be the computed candidate, $a=\min p_i$, $b=\min q_i$, and
$H'$ the order-free bound on residual weights. Every excluded tree has
$I\ge I_*+k+1$, so it costs at least
\[
 B_k=(a+b)(I_*+k+1)+(a+2b)n+H'.
\]
Thus $C_k\le B_k$ proves global optimality. Sorting and the two-queue
bound cost $O(n\log n)$ and are included in the counters. The constructor
uses $k=2$ and attempts the shortcut only when
$s\le\operatorname{bit\_length}(n)$. Its attempted work is $O(n\log n)$;
a failed certificate or skipped size falls back to Knuth. The total worst
case remains quadratic. Arbitrary growth of $k$ is exponential, and wide
shape spaces remain quadratic even for fixed $k$.

\paragraph{Validation, failed precursor, and prior work.}
The implementation passed $1{,}200$ restricted-optimum comparisons against
all shapes through eight keys with budgets $0,\ldots,4$, and $3{,}000$
full mixed-weight checks against cubic DP. It separately checks direct
tree costs, both residual-bound implementations, and exact fallback.
An additional $104$ family checks include zeros, geometric weights,
spikes, and key-only/gap-only controls. The seeded size sweep is retained.
A precursor substituting ordered depth bounds into the old certificate
passed $480$ exhaustive checks but rescued none of $24$ near-uniform
fixtures at its $32$-iteration budget. This negative experiment did not
prove optimal multiplier search useless; its cost and the failed
perfect-tree hypothesis motivated expanding the candidate class.

On the $1{,}023$-key periodic fixture with weights $16$ and $17$, the
excess-two DP finds cost $312{,}154$ and excludes other shapes with bound
$312{,}185$. It uses $27{,}570$ states, $137{,}835$ transition checks,
$6{,}052$ sorting comparisons, and $1{,}022$ merge comparisons, versus
Knuth's $523{,}776$ states and $1{,}046{,}532$ root candidates.
At $n=1{,}024$, a seeded amplitude-four perturbation uses $48{,}003$ states
and $240{,}000$ transition checks, plus $13{,}784$ sorting and $1{,}264$
merge comparisons. It certifies cost $19{,}477{,}996$.

Six matched trials per fixture use all orderings of Knuth, the preceding
baseline-price hybrid, and the excess-two constructor. Median excess-two
runtime divided by Knuth runtime was $0.23600$ and $0.44831$ on
amplitude-four perturbations at $n=1{,}023$ and $1{,}024$, and $0.20790$
and $0.17101$ on the periodic fixture at $n=1{,}023$ and $2{,}047$.
Each certified in all six trials. Ordinary random input at $n=1{,}023$
failed and cost $1.34949$ times Knuth; amplitude-sixteen input at
$n=2{,}047$ failed and cost $1.13631$ times Knuth. The skipped-width
control at $n=768$ ran the same Knuth routine and measured $0.93148$,
illustrating noise rather than a computational gain. These are exploratory
M4/Python timings including all preprocessing and reconstruction. Periodic
fixtures are repeated identically; random and perturbation inputs vary
between trials. All $42$ observations are retained.

The dyadic/prefix-code viewpoint and sequential frontier DP are established
ideas. Becker~\cite{becker} gives a sequential mixed-weight DP for a
near-minimal \emph{height} restriction in quadratic time. The band above
uses excess \emph{total unweighted path length}, a different restriction;
it does not improve Becker's stated problem. Priority of this particular
potential, band bound, and exclusion combination remains unestablished.

\section{A split-height construction covering every input size}
The preceding shape band can be wide even at zero excess. We therefore
implement the split-height/Monge mechanism already sketched by Karpinski,
Larmore, and Rytter in Section~5 of their longer preprint~\cite{klrheight}.
That section explicitly supplies the fixed-height Monge property and the
two-part construction. It is different from Section~5 of the six-page
conference version~\cite{klr}. No priority claim is made for this mechanism.
The present experiment makes the tiling, matrix padding, reconstruction,
and separate unrestricted certificate explicit.

\paragraph{Two scales of dyadic intervals.}
Fix an external height limit $H=L+U$ and write $K=2^L$, $B=2^U$.
Precompute all relative costs $C_h(i,j)$ for $0\le h\le L$ and
$0\le j-i\le2^h-1$, using
\[
 C_h(i,j)=W(i,j)+\min_r\{C_{h-1}(i,r)+C_{h-1}(r+1,j)\}.
\]
Empty intervals have zero cost; the root range enforces both child height
capacities. Monge inequalities at fixed height give monotone minimizing
roots as the endpoints move. Choose leftmost ties and bound each root by
the adjacent shorter intervals. Search ranges telescope to $O(n)$ for
each interval length. Summing over lengths and heights costs $O(n2^L)$.

Divide the unit dyadic interval into $B$ top-grid blocks. Every tree of
height at most $H$ consists of lower subtrees contained in individual
blocks, together with shallower gaps occupying aligned groups of blocks.
Let $D_b(i)$ be the minimum partial cost after consuming $i$ gaps and
reaching grid boundary $b$. Keys through $i-2$ have already been charged.
For $b>0$, put
\[
 R_b(i)=D_b(i)+p_{i-1}(U-v_2(b));\qquad R_0(0)=0.
\]
The boundary key has exactly this comparison length. There are two types
of transition:
\begin{align*}
 D_{b+1}(j)&\gets\min\{D_{b+1}(j),
 R_b(i)+C_L(i,j-1)+U W(i,j-1)\}, && i<j\le i+K,\\
 D_{b+w}(i+1)&\gets\min\{D_{b+w}(i+1),
 R_b(i)+q_i(U-\log_2 w)\}. &&
\end{align*}
The second transition requires a power of two $w\ge2$, $w\mid b$, and
$b+w\le B$. Width-one gaps are zero-key lower subtrees in the first
transition. Start at $D_0(0)=0$ and finish at $D_B(n+1)$. Dyadic alignment
proves that these paths represent every height-feasible tree and only
such trees. Backpointers recover every lower tree and every boundary key.

\paragraph{Linear matrix searches with feasible padding.}
At fixed $b$, the first transition is a banded Monge matrix in $i,j$;
row additions and the separable interval mass preserve the inequality.
Reachable counts $i$ are contiguous: the coarsest tiling of the prefix
has $\operatorname{popcount}(b)$ leaves, and successive refinement gives
every count through $bK$, truncated to $n$ before another transition.
Search output columns from $i_{\min}+1$ through
$\min(n+1,i_{\max}+K)$, each of which has a feasible input row.

SMAWK requires care outside the band. Let
$S=(H+1)(\sum p+\sum q)+1$, larger than every feasible prefix cost.
In the transposed matrix with row $j$ and candidate column $i$, use
\[
 A(j,i)=\begin{cases}
 S+j-K-i,&i<j-K,\\
 R_b(i)+C_L(i,j-1)+UW(i,j-1),&j-K\le i<j,\\
 S+i-j+1,&i\ge j.
 \end{cases}
\]
This padding preserves total monotonicity. For $i_1<i_2$, once $i_2$
is preferred at a row, it remains preferred as $j$ increases: within the
feasible band this is Monge; outside it, distances decrease toward the
band and increase away from it, and finite entries beat all padding.
When the candidates straddle opposite invalid sides, their cost difference
is $i_1+i_2+K+1-2j$, also decreasing. Ties preserve the same non-strict
implication. Thus SMAWK finds leftmost minima in $O(n)$ queries per grid
boundary. Padding by tied infinities does not supply this argument.

Across all boundaries there are $B-1$ aligned choices of coarse width,
so the second transitions and scans also cost $O(nB)$. Choosing
$L=\min(H,\lfloor H/2\rfloor+1)$ gives
\[
 O\bigl(n(2^L+2^U)\bigr)=O(n2^{H/2})
\]
work and storage including backpointers. Consequently a fixed extra height
$H=\lceil\log_2(n+1)\rceil+c$ gives $O(n^{3/2})$ attempted work at
\emph{every} size. The prototype uses sparse dictionaries; the bounded
integer state coordinates also admit direct array indexing.

\begin{lemma}[Minimum path length with a tall path]
For $1\le R\le n$, the minimum internal path length $J(n,R)$ among
$n$-key trees of external height at least $R$ is obtained by forcing a
path of $R$ internal nodes and filling the remaining available positions
in increasing order of depth. The excluded class is empty if $R>n$.
\end{lemma}
\begin{proof}
Every tree in the class contains such a path, contributing $R(R-1)/2$.
The remaining nodes form an ancestor-closed subset of its free forest:
one side root at each depth $1,\ldots,R-1$, and two roots at depth $R$.
The $n-R$ shallowest positions in this forest are ancestor-closed and
minimize the depth sum among all subsets of that size. Their addition
constructs a feasible tree attaining the bound. Directions of the forced
path do not change the forest's depth counts.
\end{proof}
A level-count implementation computes $J$ in $O(R+\log(n+1))$ work.

\begin{theorem}[Certificate against taller trees]
Let $C_H$ be the exact height-restricted optimum above, and let $Q$ be
the order-free mixed-node Huffman lower bound for residual weights
$p-a,q-b$, where $a=\min p$, $b=\min q$. If $H<n$ and
\[
 C_H\le(a+b)J(n,H+1)+(a+2b)n+Q,
\]
the constructed tree is globally optimal. If $H\ge n$, there is no
height restriction and certification is automatic.
\end{theorem}
\begin{proof}
Every excluded tree has internal path length at least $J(n,H+1)$.
Its baseline contribution is $(a+b)I+(a+2b)n$ and its residual cost
is at least $Q$. Within the height class, $C_H$ is optimal.
\end{proof}
Sorting and the residual oracle add $O(n\log n)$ work. Failed tests invoke
Knuth, so the full constructor remains quadratic in the worst case.
The certificate never assumes success from near-uniform weights.

\paragraph{Exact verification and operation evidence.}
The harness passed 450 exhaustive weighted height-class checks through
eight keys, 44 exhaustive tall-path minima, 2,000 cubic-checked mixed inputs
with complete fallback (917 successful certificates), and 104 additional
adversarial-family inputs. An independent all-root height recurrence checked
520 candidates and 261,482 stored short-table entries; the same 520 cases
also compared SMAWK with a retained divide-and-conquer matrix search.
The larger sweep has 36 fixtures at three height budgets each. Every
near-uniform fixture certified, including previously skipped sizes
384, 768, 1,000, and 1,536. All random, geometric, and gap-spike fixtures
failed the certificate, sometimes despite an optimal restricted candidate.

At $n=1{,}000$, the minimum-height candidate and full optimum are both
$18{,}973{,}413$, below the exclusion bound $18{,}974{,}193$. The optimizer
uses 117,453 short states, 220,483 short root candidates, 84,937 genuine
matrix queries, 10,268 padding queries, 16,016 prefix scans, and 5,630 coarse
transitions. The bound adds 10,648 sorting and 1,335 merge comparisons,
plus 1,000 contractions; all keys are reconstructed. Knuth uses 500,500
states and 1,002,162 root candidates. These distinct operation types are
reported separately, without treating them as equal-cost instructions.

\paragraph{Matched timing and failure controls.}
Six trials per fixture use every ordering of Knuth, the minimum-height
constructor ($c=0$), and one extra level ($c=1$). Preprocessing and
reconstruction are included; every result is checked against Knuth.
Inputs vary between trials except the repeated periodic fixture. On the
sponsor's M4 with Python~3.9.6, median per-trial runtime ratios were:
\begin{center}\small
\begin{tabular}{lrrrr}
\toprule
Fixture & $c=0$ ratio & certified & $c=1$ ratio & certified\\
\midrule
384, near-uniform & 1.22484 & 6/6 & 2.29965 & 6/6\\
768, near-uniform & 0.94444 & 6/6 & 1.40863 & 6/6\\
1,000, near-uniform & 0.71530 & 6/6 & 1.03015 & 6/6\\
1,536, near-uniform & 0.65826 & 6/6 & 1.03916 & 6/6\\
3,072, near-uniform & 1.46288 & 0/6 & 0.70837 & 6/6\\
1,536, amplitude four & 1.67333 & 0/6 & 1.18924 & 6/6\\
1,536, ordinary random & 1.55378 & 0/6 & 1.93841 & 0/6\\
1,023, periodic $A=16$ & 1.69541 & 0/6 & 1.06171 & 6/6\\
\bottomrule
\end{tabular}
\end{center}
All 48 observations and operation counters are retained. These are mixed
practical results: small successful cases can be slower, and random cases
pay substantial overhead before quadratic fallback. At 3,072 near-uniform
keys, the minimum-height candidates are optimal but fail certification;
one extra level strengthens the exclusion bound sufficiently. For the
periodic fixture, minimum height actually excludes the true optimum.
The earlier bounded-excess solver is more efficient on that narrow class.
No broad practical superiority or improvement over the known width-limited
Larmore/KLR constructor is established by this experiment.

\section{A variable interval frontier for exact full construction}
The height test can fail even when its candidate is optimal. We now use a
class guaranteed in advance to contain an unrestricted optimum. The
short-interval/top-skeleton architecture is established work of Larmore
and Karpinski--Larmore--Rytter~\cite{larmore,klr}; the experiment compares
a uniform width envelope with a variable interval contour. Both use the
same safe mass guarantee and matrix infrastructure. This is not a priority
claim for the architecture or its restricted subquadratic regime.

\paragraph{A failed transfer of the sharper mass constant.}
The initial implementation used $2W/F_{d+3}$ as the maximum internal
subtree mass at edge depth $d$, following a literal transfer of the
Fibonacci formula in~\cite{hlm}. Its first full-input check failed: the
resulting class could be empty. The positive two-key example
\[
 p=(1,1),\qquad q=(1,100,1)
\]
explains the failure. Both trees cost 206; the total mass is 104 and the
internal child has mass 102. Thus $102>2\cdot104/3$ at depth one. Charging
an additional comparison at every gap is a tree-independent constant and
does not remove this violation. We preserve the counterexample and use
the following independently justified, weaker constant in our model.
This does not refute the KLR architecture or its asymptotic guarantees.

\begin{lemma}[A sufficient mixed-weight subtree bound]\label{lem:mass}
If $S$ is an internal subtree at edge depth $d$ in an optimal full OBST
of total outcome mass $W$, then
\[
 W(S)\le W/F_{d+1},\qquad F_1=F_2=1.
\]
\end{lemma}
\begin{proof}
Let $X$ be an internal grandchild of $v$, with parent $y$, and let $O$
be the mass inside $v$'s subtree but outside $y$'s, including key $v$.
For a straight two-edge path, rotating $y$ upward changes cost by
$O-p_y-W(X)$. For a bent path, promoting $X$ with a double rotation
changes cost by $O-p_X-W(X)$. Optimality therefore implies $O\ge W(X)$
in either case. These formulas include all external gap weights.

Going upward from $S$, write $T_0=S,T_1,\ldots,T_d=T$. Nonnegativity
gives $W(T_1)\ge W(T_0)$, and the grandchild inequality gives
$W(T_i)\ge W(T_{i-1})+W(T_{i-2})$ for $i\ge2$.
Induction gives $W(T_d)\ge F_{d+1}W(S)$. Weak inequalities include
zero weights and ties.
\end{proof}
This lemma concerns internal subtrees. A heavy interior external gap can
have nearly all the mass at depth two; zero-key subtrees remain permitted.
The Fibonacci approach is prior work~\cite{hlm}; sharpness of this
particular sufficient bound is not claimed.

\paragraph{Store a contour instead of a global width.}
For fixed $d$, let $E_d(i)$ be the largest $j>i$ satisfying
$F_{d+1}W(i,j)\le W$, or $i$ if there is none. Empty intervals always
have stored cost zero. Nonnegative weights imply that $E_d(i)$ is
nondecreasing and the domain is closed under taking subintervals.
A two-pointer scan computes it in $O(n)$ work. Write
\[
 S_d=\sum_{i=0}^n(E_d(i)-i).
\]
By Lemma~\ref{lem:mass}, every internal subtree at depth $d$ of an
optimal tree belongs to the domain. For comparison, the uniform version
sets $\ell_d=\max_i(E_d(i)-i)$ and stores
\[
 U_d=\ell_d(n+1)-\ell_d(\ell_d+1)/2\ \ge S_d
\]
intervals. Both classes contain an unrestricted optimum.

\begin{lemma}[Root-search work on a hereditary contour]
Computing all unrestricted interval optima in this domain takes
$O(n+S_d)$ work, with at most $3S_d$ tested roots.
\end{lemma}
\begin{proof}
Bucket only existing cells by length in $O(n+S_d)$ work. Knuth's
recurrence applies because every child and adjacent shorter interval is
present. Let $r(i,j)$ be the leftmost optimal root and put
$t(i,j)=r(i,j)-i\ge0$. For length at least two, the root count is
$2+t(i+1,j)-t(i,j-1)$. Sum over the domain. A positive occurrence
$t(a,b)$ cancels its negative occurrence whenever $(a,b+1)$ is present.
An unmatched positive occurrence requires $b=E_d(a)$ and
$E_d(a-1)=E_d(a)$. There is at most one per row, of size at most
$E_d(a)-a$. Discard unmatched negative occurrences. Adding the length-one
cells proves the $3S_d$ bound. Scanning nonexistent diagonal cells would
lose this saving; the buckets are essential.
\end{proof}

\paragraph{Compose the exact interval costs.}
Use $2^d$ top-grid blocks and the preceding dyadic transitions, replacing
height-limited lower costs by these unrestricted interval costs. A block
starting after $i$ gaps can finish after $j$ exactly when
$i<j\le E_d(i)+1$. It contributes $C(i,j-1)+dW(i,j-1)$ and the
starting boundary key cost $p_{i-1}(d-v_2(b))$ at $b>0$.
Coarse gaps retain their aligned power-of-two skip transitions. Every
optimal tree is represented, while every reconstructed candidate is a
valid full tree. Thus the resulting optimum is unrestricted.

The variable band needs another implementation detail. Reachable input
rows, and their union of feasible output columns, may have holes. Scan
once per top boundary to retain exactly output columns covered by some
reachable row. Put $L(j)=\min\{i:E_d(i)+1\ge j\}$, a nondecreasing
sequence with $L(j)\le j-1$. In the transposed matrix, use finite true
entries when $L(j)\le i<j$, and padding
\[
 Z+L(j)-i\quad(i<L(j)),\qquad
 Z+i-j+1\quad(i\ge j),\qquad Z=(n+d+1)W+1.
\]
All finite prefix costs are smaller than $Z$. Within the band, exact
interval costs are Monge. Across padding, opposite-side differences
\mbox{$i_1+i_2+1-j-L(j)$} are nonincreasing; same-side preferences also
move right. Thus the matrix is totally monotone,
even after removing arbitrary rows and columns, and every searched column
has a feasible minimum. SMAWK, scans, and coarse transitions together
cost $O(n2^d)$. Backpointers recover every key and gap contribution.

\begin{theorem}[Exact mass-frontier constructor]
The adaptive implementation constructs a full optimal BST in
\[
 O\!\left(n\log(n+1)+\min_{d\ge0}\{S_d+n2^d\}\right)
\]
arithmetic work. The uniform-envelope implementation replaces $S_d$ by
$U_d$. At the selected depth, storage is $O(n+S_d+n2^d)$ or
$O(n+U_d+n2^d)$ respectively. Both remain quadratic in the worst case.
\end{theorem}
\begin{proof}
The interval and composition arguments establish exactness at every depth.
Start with depth zero and the full Knuth table. Evaluate successive depths
while $n2^d$ is smaller than the best work proxy so far; each contour costs
$O(n)$. No later depth can improve the proxy because its top term alone
is too large. At most $O(\log(n+1))$ depths are evaluated. Depth zero
uses Knuth directly, omitting a redundant top pass and its lower-order
$n$ proxy term. The preceding work and storage bounds now apply.
\end{proof}
This constructor has no posterior certificate failure followed by a second
solve. It may select the full table initially. The Python prototype uses
dictionaries; bounded integer top coordinates also allow direct arrays.
When each key has weight at least $a>0$ and $W=O(na)$, we have
$S_d=O(n^2/\phi^d)$, where $\phi=(1+\sqrt5)/2$. Balancing against
$n2^d$ recovers the known exponent $1+\log_{2\phi}2\simeq1.59023$.

\paragraph{A scoped asymptotic separation.}
Let $n$ range over fourth powers and $L=n^{3/4}$. Give every key and gap
weight $A=n^2$, except a central block of $L$ keys and its $L+1$ gaps,
all of weight one. Every weight is positive and $W=\Theta(nA)$.
An admitted interval contains at most $m=O(n/F_{d+1})$ heavy keys.
Intervals wholly inside the light block number $O(L^2)$. Other starts
contribute $O(nm)$ ordinary endpoints; only $O(m)$ outside starts can
reach across a block boundary and gain $O(L)$ endpoints, and starts
inside add at most $O(Lm)$ heavy endpoints. Consequently
\[
 S_d=O(n^2/F_{d+1}+L^2).
\]
With $d=\lfloor\log_{2\phi}n\rfloor$, the adaptive work is
$O(n^{1+\log_{2\phi}2}+n^{3/2})=O(n^{1.59024})$.

Every uniform-envelope plan, however, has proxy $\Omega(nL)$.
If $2^d\ge L$, the top term suffices. Otherwise
$F_{d+1}=O(L^{\log_2\phi})$, so $W/F_{d+1}$ exceeds the entire light
block's $O(L)$ mass. Hence $\ell_d\ge L$ and the uniform table has
at least $nL/2$ cells. This gives $\Omega(n^{1.75})$ at every depth.
The separation is within the specified architecture, not against every
prior OBST algorithm or combinations of our shortcuts. In particular,
this light block also offers an opportunity for certified contraction.

Both modes were checked against Knuth on this positive family at
$n=16,81,256,625$, with cubic checks also at the first two sizes. Larger
instances were used only for plan counting. At $n=65{,}536$, the chosen
adaptive and uniform proxies were 119,041,540 and 361,608,522, both at
depth nine. These are $S_d+n2^d$ and $U_d+n2^d$, not measured full-solve
operation totals or timings. All parameter counts are retained.

\paragraph{Independent verification.}
The harness passed 960 weighted exhaustive-shape comparisons through
eight keys on arbitrary nondecreasing contours, including infeasible
classes. It checked 914 interval costs independently and 90,193 top
minima by brute search. Another 2,400 full solves, using both modes on
1,200 inputs through 40 keys, matched cubic DP and reconstructed costs.
All 24,600 internal subtrees in the cubic-optimal trees satisfied
Lemma~\ref{lem:mass}. Families include zeros, ties, geometric variants,
spikes, and key-only and gap-only controls. A further 56 fixtures through
2,048 keys matched Knuth in both modes, including ordinary random inputs
that failed the preceding height certificates.

At 2,048 random keys, the adaptive version chose depth six and returned
cost 21,053,524. It used 309,261 interval cells and 616,642 root candidates,
versus Knuth's 2,098,176 cells and 4,215,527 candidates. Additional work
included 24,035 contour mass checks across seven parameter depths,
1,245,431 genuine matrix queries, 149,319 padding queries, 131,136 scans
of each row/column type, and 107,859 coarse transitions. The uniform
version stored 337,550 cells. On log-uniform weights the corresponding
cell counts were 308,487 versus 431,951; on a single heavy key they were
1,048,576 versus 1,573,376. Geometric inputs selected the full table.
All operation categories and parameter levels are retained separately.

\paragraph{Matched timings.}
All six orders of Knuth, uniform width, and adaptive contour are used per
fixture, for 48 trials. Preprocessing and reconstruction are timed; every
result is directly evaluated and checked against Knuth. Inputs vary except
for deterministic geometric and spike fixtures. The timing light window
uses weights near 1,000 outside a block of $\lfloor n^{3/4}\rfloor$ keys
and weight one inside, with varying position; it is separate from the
$A=n^2$ asymptotic family. M4/Python~3.9.6 median paired ratios were:
\begin{center}\small
\begin{tabular}{lrrr}
\toprule
Fixture & uniform/Knuth & adaptive/Knuth & adaptive/uniform\\
\midrule
512, random & 1.59908 & 1.61204 & 0.96602\\
1,024, random & 1.56592 & 1.55650 & 1.00412\\
2,048, random & 0.90002 & 0.91021 & 1.00096\\
2,048, light window & 1.13693 & 0.77546 & 0.65523\\
2,048, near-uniform & 0.87926 & 0.91295 & 1.02587\\
1,024, log-uniform & 1.67408 & 1.47983 & 0.88390\\
512, geometric & 0.98060 & 1.04487 & 1.01772\\
2,048, heavy key & 1.50542 & 1.04249 & 0.70113\\
\bottomrule
\end{tabular}
\end{center}
The smaller random cases remain slower than Knuth. The contour's saving
over uniform width is clearest on localized structure; ordinary random
inputs show little timing difference. Both geometric variants execute the
same Knuth kernel after selection, so the sub-unit ratio there is noise.
These instrumented Python measurements do not establish production-language
or broad practical superiority. Every raw observation is retained.

\section{Exact checks and preliminary measurements}
The code uses nonnegative integer weights, an independent cubic recurrence,
Knuth's recurrence, and direct evaluation of reconstructed trees.
The basic harness additionally enumerates all tree shapes for small sizes.
CPU experiments ran on the sponsor's M4 Mac mini. The reported quantities
are algorithmic counts, not inferred asymptotic bounds from elapsed time.

\begin{itemize}
\item Single-key contraction: 22,143 exhaustive weight assignments in
$\{0,1,2\}$ through $n=4$, and 6,000 seeded random instances through $n=20$.
Every cost and reconstructed tree matched the cubic oracle.
\item Block theorem: 19,486 qualifying interval contractions across 4,500
seeded instances through $n=16$ matched independent cubic costs.
\item Block algorithm: 3,000 random instances through $n=20$, including
fixed and adaptive thresholds, matched the cubic oracle and tree evaluation.
Additional small cases check threshold selection against every threshold.
\end{itemize}

Table~\ref{tab:states} records one seeded instance per family at $n=512$.
The baseline has 131,328 states. The block column includes all states
inside selected blocks, not just the core. Results from different solvers
use their separately recorded seeds and are not paired timing experiments.
\begin{table}[ht]
\centering
\caption{Preliminary exact DP-state counts, including block solves.}
\label{tab:states}
\begin{tabular}{lrr}
\toprule
Family & Single-key pass + DP & Adaptive blocks + DP\\
\midrule
Uniform mixed weights & 131,328 & 131,328\\
Near-uniform mixed weights & 131,328 & 131,328\\
Random mixed weights & 128,778 & 130,817\\
Log-uniform weights & 116,886 & 94,964\\
Keys only (control) & 92,235 & 88,508\\
Sparse heavy separators & --- & 10,344\\
\bottomrule
\end{tabular}
\end{table}

In the heavy-separator instance all gaps and ordinary keys have integer
weights in $\{1,2,3\}$, and 22 randomly selected keys have weight $10^6$.
The block algorithm retains exactly 22 keys. Its state reduction is about
12.7-fold, but this is not a claim of a 12.7-fold runtime speedup.
On ordinary random mixed weights the gain is small.
The single-key pass contracts every key for the tested sequence
$p_i=4^i,q_i=4^i$, but not for most of the ratio-two sequence.
Geometric integer inputs also illustrate why bit complexity must be
distinguished from arithmetic-operation complexity.

\section{A useful relaxation that failed as a general solver}\label{sec:unordered}
Allow arbitrary assignment of key weights to internal positions and gap
weights to external positions. In this relaxed problem, repeatedly take
the two smallest current external weights $a,b$ and the smallest remaining
internal weight $p$, replace them by $a+b+p$, and add that sum to the cost.

To prove the rule, take a deepest internal node of an optimum.
Its children are leaves. Swapping the smallest internal weight into it
and the two smallest external weights into its children cannot increase
cost. Contract that cherry and induct. Sorting and a heap give
$O(n\log n)$ time. This elementary extension of the Huffman exchange
argument is not claimed as novel.

The relaxed cost is a valid lower bound for full OBST. Equality with a
feasible ordered tree certifies optimality. It agreed with an independent
shape-enumeration and rearrangement oracle on 22,143 exhaustive and 2,700
random instances in the basic harness.

However, using the relaxation on demand in an exact sparse DP was poor
on random inputs. At $n=128$ it visited roughly 4,300--4,600 exact intervals
instead of Knuth's 8,256, but spent approximately 3.9--4.1 million weight
comparisons constructing bounds and inspected over 100,000 candidate roots.
Fewer exact states did not mean less computation.

Nor does the relaxation's interval matrix inherit the required quadrangle
inequality. For $p=(9,1,12)$, $q=(5,1,18,10)$, its costs $H$ satisfy
\[
 H(1,2)+H(2,3)=60>54=H(1,3)+H(2,2).
\]
Thus applying Knuth-style matrix optimization directly to this bound
would be unjustified.

Two additional exact counterexamples exclude tempting local certificates.
With $p=(3,5,18)$, $q=(13,14,4,8)$, the tree with root 2 and left subtree
$0(\varnothing,1)$ costs 127, every single-rotation neighbor costs at least 134,
yet the balanced tree costs 125.
With $p=(18,8,18)$, $q=(17,11,1,13)$, the balanced tree costs 164 and every
distinct tree obtained by promoting a descendant to the root of an existing
subtree costs at least 169; the tree $0(\varnothing,2(1,\varnothing))$ costs 157.
These are barriers to those local certificates, not complexity lower bounds.

\section{An ordered depth-chain lower bound}
Let $l_i$ be the comparison length of key $i$ and $g_i$ the depth of gap $i$.
For every valid BST, adjacent keys have different lengths and
\[
 g_0=l_0,\quad g_n=l_{n-1},\quad
 g_i=\max(l_{i-1},l_i)\quad(0<i<n).
\]
This identity is existing work; see Becker~\cite{becker}, Lemma~1.
Let $K_d$ count keys of length $d$ and $G_d$ count gaps at depth $d$.
The root and child-count identities are
\[
 K_1=1,\qquad G_d+K_{d+1}-2K_d=0\quad(1\le d\le n),
 \qquad K_{n+1}=0.
\]

Permit any sequence in $\{1,\ldots,n\}^n$ with unequal adjacent entries.
For arbitrary multipliers $\lambda_0,\ldots,\lambda_n$, minimize
\begin{equation}\label{eq:dual}
 \sum_i p_i l_i+\sum_i q_i g_i
 +\lambda_0(K_1-1)
 +\sum_{d=1}^n\lambda_d(G_d+K_{d+1}-2K_d).
\end{equation}
Call the minimum $L(\lambda)$. Every actual tree is permitted and has zero
penalty, so $L(\lambda)\le\OPT$ for every multiplier vector.
The expression is a chain with key energy
$U_i(b)=p_i b+\lambda_{b-1}-2\lambda_b$, gap energy
$V_i(b)=q_i b+\lambda_b$, and constant $-\lambda_0$.
For an interior transition into length $b$, take the smaller of
\[
 \min_{a<b}F(a)+V_i(b),\qquad
 \min_{a>b}\{F(a)+V_i(a)\}.
\]
Prefix and suffix minima give $O(H)$ work per key for a domain of $H$ lengths.
Using $H=n$ would cost quadratic time even for one bound evaluation.

\begin{theorem}[Exact affine-tail compression]
Fix $h\ge1$, arbitrary $\lambda_0,\ldots,\lambda_h$, and a slope
$0\le s\le\min_i p_i$. Extend
$\lambda_d=\lambda_h+s(d-h)$ for $d>h$.
The minimum in~\eqref{eq:dual} is unchanged if lengths are restricted to
$1,\ldots,\min(n,h+3)$.
\end{theorem}
\begin{proof}
For $b\ge h+1$, the unary key energy has slope $p_i-s\ge0$ and the gap
energy has slope $q_i+s\ge0$. Both are nondecreasing in this tail.
If some $l_i>h+3$, replace it by a member of
$\{h+1,h+2,h+3\}$ unequal to its current neighbors. At most two values
are forbidden. The key energy does not increase. Each adjacent maximum
is either unchanged or decreases within the affine tail, so neither gap
energy increases. Endpoint gaps obey the same argument. Repeat for all
large lengths. The restricted domain is a subset of the original domain,
and every original sequence has a no-more-expensive restricted sequence;
therefore the minima are equal.
\end{proof}

Taking $h=\lceil\log_2(n+1)\rceil+2$ gives an $O(n\log n)$ oracle.
This is a statement about minimization of the relaxed energy, with no height
assumption on the unknown optimal BST. Optimizing only the head multipliers
requires folding all tail constraint violations into the derivative for
$\lambda_h$. The implementation uses exact integer multipliers and costs
scaled by 1024. Its subgradient ascent is heuristic; the bound is valid
at every iteration regardless of convergence.

A stronger oracle requires exactly one length1 key. For each candidate
root, the left and right chains then contain only lengths at least2.
One forward and one reverse pass supply their closing costs for every
root position, giving all forced-root bounds in $O(nH)$ total time.
The projection above preserves the unique root, so tail compression remains
valid for each marginal. Retaining the maximum bound seen for each root,
then taking the minimum over roots, is also a valid lower bound.

Both oracles passed 50 comparisons against exhaustive length-sequence
minimization, 5,000 arbitrary-multiplier bound checks against cubic OBST,
and 1,050 direct comparisons of the full and shortened oracles.
Forced-root marginals were additionally checked against exact
root-constrained costs.

\paragraph{A structural gap.}
The sequence $l=(3,1,2,3,2)$ has
$g=(3,3,2,3,3,2)$, a unique root, and all the level-count equalities.
It is not a valid BST depth sequence. For
\[
 p=(3,27,21,16,12),\qquad q=(18,22,27,22,12,22),
\]
it costs 470, while the true optimum is 472. Its penalty is zero for every
multiplier vector. Hence neither oracle family can ever exceed 470 on
this instance. More optimization iterations cannot remove this obstruction.

\paragraph{Preliminary strength and cost.}
At $n=256$, normalized gaps $(\OPT-\mathrm{LB})/W$ for the chain bound
versus order-free Huffman were 0.15003 versus 0.22080 on a random mixed
instance, and 0.00873 versus 0.49294 on an alternating instance.
However, 300 ascent iterations used 1,082,368 key-depth cells including
initialization, compared with 32,896 interval states for Knuth.
The unique-root variant uses two passes and still leaves 254 of 256
potentially improving roots on the random instance, and 255 on the
alternating instance. Only two survive on the gap-spike instance.
These results identify a lower-bound component and its limitations;
they do not establish a faster general exact solver.

\section{Exact depth ordering and an arbitrary-window obstruction}
Call a positive integer sequence \emph{ranked} if every nonempty contiguous
interval has a unique minimum. Equivalently, equal labels have a smaller
label between them. This is vertex ranking on a path with the label order
reversed; vertex ranking and bounded-path variants are existing
notions~\cite{ranking,localranking}. The following results explain the
ordering information absent from the preceding relaxation; no novelty claim
is made for this characterization.

\begin{theorem}[Depth characterization]
A sequence $l\in\{1,\ldots,n\}^n$, $n\ge1$, is the inorder
key-comparison-length sequence of a BST if and only if it is ranked,
$K_1=1$, and every level-flow equality holds, with gaps derived from $l$
as above.
\end{theorem}
\begin{proof}
Necessity follows from the unique shallowest key in each inorder interval
and the two children of every internal node. For sufficiency, delete all
labels at most $d$. Let $U_d$ count the remaining nonempty runs, and let
$A_d$ count the deleted labels. If the deleted labels form $R$ runs and
include $e$ endpoints, then the number of gaps of length at most $d$ is
$A_d-R+e$, whereas $U_d=R+1-e$. Thus
\[
 G_{\le d}=A_d+1-U_d.
\]
Summing the level-flow equalities from 1 through $d$, and using $K_1=1$,
also gives
\[
 G_{\le d}=A_d+1-K_{d+1}.
\]
Hence $U_d=K_{d+1}$. Ranking permits at most one label $d+1$ per remaining
run, so equality forces exactly one. Starting with the unique length1
root, every remaining nonempty interval therefore receives its root at
the next level. This constructs the tree with exactly the supplied lengths.
\end{proof}

Ranking has a linear recognizer: scan with an increasing stack, pop larger
labels, reject an equal top, and otherwise push. A rejected pair has no
smaller intervening label; conversely every such pair causes rejection.
Each entry is pushed once and popped at most once. This establishes cheap
recognition, not cheap minimization over all valid sequences.

\begin{theorem}[Bounded-window obstruction]
For every integer $w\ge2$ and $t\ge2$, there is a strictly positive full-OBST
instance with $n=w+3$ keys and an inadmissible depth sequence satisfying
all level-flow equalities, a unique length1 root, and ranking on every
interval of at most $w$ keys. Its cost is at most
$(7/8+1/(8t))\OPT$. Consequently every lower bound obtained by pricing
the level-flow equalities over this window-constrained domain has the same
upper limitation, regardless of the multipliers.
\end{theorem}
\begin{proof}
Take
\[
 l=(3,1,2,3,4,\ldots,w+1,2).
\]
The ascending middle portion has $w-1$ entries. Its two enclosing length2
labels form the only interval with a repeated minimum, of width $w+1$.
Every shorter window passes. Direct substitution into the gap formula
verifies $K_1=1$ and every level-flow equality. The sequence is invalid:
the unique root is key1, but its singleton left subtree has length3,
and two keys to its right both claim length2.

Give every gap weight1 and every key weight1 except
\[
 p_1=3B,\qquad p_2=p_{n-1}=B.
\]
The unit-weight outcomes contribute
$E=w^2+4w+10$ in the inadmissible sequence, whose cost is $7B+E$.
Every actual tree costs at least $8B$. If key1 is the root, the two
$B$-weight keys lie to its right and their comparison lengths sum to at
least5, giving $3B+5B$. Otherwise key1 has length at least2 and each
$B$-weight key has length at least1, giving $6B+2B$.

Set $B=tE$. The inadmissible cost satisfies
\[
 7B+E=\left(\frac78+\frac{1}{8t}\right)8B
       \le\left(\frac78+\frac{1}{8t}\right)\OPT.
\]
Its level-flow penalty is zero for every multiplier vector, so every
priced minimum is at most this cost.
\end{proof}

For $t=16$ the upper ratio is $113/128$, independently of $w$; as $t$
increases it approaches $7/8$. This rules out universal exactness of the
specified window-minimum repair, not other local constraints, adaptive
nonlocal cuts, or subquadratic full-OBST algorithms.
Equivalently, at every $n\ge5$ the construction passes all windows of up
to $n-3$ keys.

\paragraph{Burying the obstruction below a forced prefix.}
Let a current instance have mass $W$ and a feasible tree of cost $U$.
Prepend a key of weight $A=W+U+2$ and a new leftmost gap of weight1.
Using the new key as root costs $A+1+W+U=2A-1$, while every other root
costs at least $2A$ from that key alone. Hence the new root is forced and
$\OPT_{\rm new}=A+1+W+\OPT_{\rm old}$.
The transformation $l\mapsto(1,l_0+1,\ldots,l_{n-1}+1)$ preserves zero
flow and local window ranking, and changes the inadmissible cost by the
same $A+1+W$. Repeating this $D$ times places the strict gap beneath the
true optimal top-$D$ right-chain prefix. Thus fixing any constant number
of top levels cannot make the remaining window/flow relaxation exact.
The constant relative gap is not claimed after this embedding, because
the forcing weights add a potentially large common cost.

\paragraph{Independent verification.}
The harness enumerated all 873,612 sequences in $\{1,\ldots,n\}^n$ for
$1\le n\le7$. The characterization accepted exactly the independently
generated BST depth sequences: respectively 1, 2, 5, 14, 42, 132, and 429.
The stack recognizer agreed with direct interval scans on every sequence.
Width3 and exact flow already admit four invalid sequences at $n=6$ and
34 at $n=7$. The positive family was checked at $t=16$ for all 127 widths
$2\le w\le128$: exact costs, window validity, zero flow, and the tree-cost
lower bound all passed. Knuth and direct tree evaluation checked every
instance; independent cubic DP also checked $w\le24$.
Another 120 cases combine $2\le w\le16$ with forced prefix depths
$1\le D\le8$. Cubic DP, Knuth, direct costs, and inspection of the optimal
prefix verify the exact additive identity and preserved strict gap.

\section{Status and reproducibility}
The general subquadratic question remains open. This checkpoint provides
exact contraction, bounded-excess and split-height constructions, variable
interval frontiers, global certificates,
and scoped obstructions to depth relaxations. Mathematical novelty, broad
practical benefit, and comparisons against other optimized prior implementations
remain to be assessed.

Standard-library Python code, seeded exact results, and matched timing
observations are under \texttt{obst/code} and \texttt{obst/results}.
The README gives reproduction commands. The active objective remains to
avoid most of the quadratic table on a substantially broader class, or
establish a rigorously scoped complexity lower bound.

\begingroup\small
\begin{thebibliography}{9}
\setlength{\itemsep}{0pt}
\raggedright
\bibitem{knuth}
D. E. Knuth. Optimum binary search trees.
\emph{Acta Informatica} 1, 14--25, 1971.
\href{https://doi.org/10.1007/BF00264289}{doi:10.1007/BF00264289}.
\bibitem{larmore}
L. L. Larmore. A subquadratic algorithm for constructing approximately
optimal binary search trees. \emph{Journal of Algorithms} 8(4),
579--591, 1987.
\href{https://doi.org/10.1016/0196-6774(87)90052-6}{doi:10.1016/0196-6774(87)90052-6}.
\bibitem{klr}
M. Karpinski, L. L. Larmore, and W. Rytter.
Sequential and parallel subquadratic work algorithms for constructing
approximately optimal binary search trees. \emph{Proceedings of SODA},
36--41, 1996.
\href{https://theory.cs.uni-bonn.de/marek/publications/seq_and_par_subquadr.pdf}{Author-hosted manuscript}.
\bibitem{klrheight}
M. Karpinski, L. L. Larmore, and W. Rytter. Same title as~\cite{klr},
longer preprint version, Section~5: Short height limited trees.
\href{https://www.researchgate.net/publication/2710579_Sequential_and_Parallel_Subquadratic_Work_Algorithms_for_Constructing_Approximately_Optimal_Binary_Search_Trees}{Archived preprint text}.
\bibitem{hlm}
D. S. Hirschberg, L. L. Larmore, and M. Molodowitch.
Subtree weight ratios for optimal binary search trees. Technical Report
86-02, University of California, Irvine, 1986.
\href{https://ics.uci.edu/~dhirschb/pubs/TR86-02.pdf}{Author-hosted report}.
\bibitem{becker}
P. Becker. Optimal binary search trees with near minimal height. 2010.
\href{https://arxiv.org/abs/1011.1337}{arXiv:1011.1337}.
\bibitem{ranking}
D. C. McDonald. On-line vertex ranking of trees.
\emph{SIAM Journal on Discrete Mathematics} 29(1), 145--156, 2015.
\href{https://arxiv.org/abs/1401.2669}{arXiv:1401.2669}.
\bibitem{localranking}
P. Bose, V. Dujmovi\'{c}, M. Javarsineh, and P. Morin.
Asymptotically optimal vertex ranking of planar graphs. 2022 version.
\href{https://arxiv.org/abs/2007.06455}{arXiv:2007.06455}.
\end{thebibliography}
\endgroup
\end{document}
